% ── Non-Hermitian Photonics — Project Macros ────────────────
% Domain-specific commands used throughout the monograph.
% Keep this file lean; add only macros used in multiple chapters.

% ── Maxwell fields and operators ────────────────────────────
\newcommand{\EE}{\mathbf{E}}             % electric field
\newcommand{\HH}{\mathbf{H}}             % magnetic field
\newcommand{\BB}{\mathbf{B}}             % magnetic flux density
\newcommand{\DD}{\mathbf{D}}             % displacement field
\newcommand{\JJ}{\mathbf{J}}             % current density
\newcommand{\FF}{\mathbf{F}}             % generic vector field
\renewcommand{\SS}{\mathbf{S}}           % Poynting vector or generic vector
\newcommand{\rr}{\mathbf{r}}             % position vector
\renewcommand{\aa}{\mathbf{a}}           % lattice vector
\newcommand{\kk}{\mathbf{k}}             % wave vector
\newcommand{\GG}{\mathbf{G}}             % reciprocal lattice vector
\newcommand{\nn}{\hat{\mathbf{n}}}       % unit normal
\newcommand{\curl}{\nabla\times}         % curl operator
\newcommand{\divg}{\nabla\cdot}          % divergence
\newcommand{\grad}{\nabla}               % gradient

% ── Permittivity and permeability ───────────────────────────
\newcommand{\eps}{\varepsilon}            % permittivity
\newcommand{\epsr}{\varepsilon_{\mathrm{r}}}  % relative permittivity
\newcommand{\mur}{\mu_{\mathrm{r}}}      % relative permeability
\newcommand{\epsc}{\tilde{\varepsilon}}   % complex permittivity (with tilde)

% ── Operators ───────────────────────────────────────────────
\newcommand{\Lop}{\hat{\mathcal{L}}}     % generic linear operator
\newcommand{\Mop}{\hat{\Theta}}          % Maxwell curl-curl operator
\newcommand{\adj}{\dagger}               % Hermitian adjoint
\newcommand{\transpose}{\mathsf{T}}      % transpose
\DeclareMathOperator{\tr}{tr}             % trace
\DeclareMathOperator{\diag}{diag}         % diagonal matrix
\DeclareMathOperator{\Res}{Res}           % residue
\renewcommand{\Im}{\operatorname{Im}}      % imaginary part
\renewcommand{\Re}{\operatorname{Re}}      % real part
\DeclareMathOperator{\sgn}{sgn}           % sign function

% ── Non-Hermitian specific ──────────────────────────────────
\newcommand{\PT}{\mathcal{PT}}            % PT symmetry
\newcommand{\Smat}{\mathbf{S}}           % scattering matrix
\newcommand{\Tmat}{\mathbf{T}}           % transfer matrix
\newcommand{\omtil}{\tilde{\omega}}       % complex frequency
\newcommand{\GBZ}{\mathrm{GBZ}}          % generalized Brillouin zone
\newcommand{\BZ}{\mathrm{BZ}}            % Brillouin zone
\newcommand{\NHSE}{\mathrm{NHSE}}        % non-Hermitian skin effect
\newcommand{\wnum}{W}                     % winding number
\newcommand{\Chern}{\mathcal{C}}          % Chern number

% ── Berry phase and topology ────────────────────────────────
\newcommand{\Acal}{\mathcal{A}}           % Berry connection
\newcommand{\Fcal}{\mathcal{F}}           % Berry curvature
\newcommand{\Qfac}{Q}                     % quality factor

% ── Biorthogonal inner product ──────────────────────────────
\newcommand{\bket}[1]{\left|#1\right\rangle}           % right eigenvector
\newcommand{\bbra}[1]{\left\langle\!\left\langle#1\right|\right.}  % left eigenvector (double angle)
\newcommand{\biinner}[2]{\left\langle\!\left\langle #1 \,\middle|\, #2 \right\rangle\!\right\rangle}  % biorthogonal inner product

% ── Commonly used abbreviations in math ─────────────────────
\newcommand{\cc}{\mathrm{c.c.}}          % complex conjugate
\newcommand{\hc}{\mathrm{h.c.}}          % Hermitian conjugate
\newcommand{\etal}{\textit{et~al.}}       % et al. in text

% ── Units and constants ─────────────────────────────────────
\newcommand{\vac}{0}                      % vacuum subscript
\newcommand{\epsvac}{\varepsilon_0}       % vacuum permittivity
\newcommand{\muvac}{\mu_0}               % vacuum permeability
\newcommand{\cvac}{c}                     % speed of light in vacuum
